\documentclass[12pt]{article}
\usepackage[utf8]{inputenc}
\usepackage{amsmath,amssymb,amsthm}
\usepackage[margin=1in]{geometry}

\newtheorem{theorem}{Theorem}
\newtheorem{lemma}[theorem]{Lemma}
\newtheorem{definition}[theorem]{Definition}

\title{Problem 153: Investigating Gaussian Integers}
\author{Project Euler Solution}
\date{}

\begin{document}
\maketitle

\section{Problem Statement}

For $1 \le n \le N = 10^8$, find $\displaystyle\sum_{n=1}^{N} R(n)$, where $R(n)$ is the sum of the real parts of all Gaussian integer divisors of $n$ with positive real part.

\section{Mathematical Development}

\begin{definition}
A \emph{Gaussian integer} is an element of $\mathbb{Z}[i] = \{a+bi : a,b \in \mathbb{Z}\}$. For $g = a+bi$, its \emph{norm} is $N(g) = a^2+b^2$.
\end{definition}

\begin{theorem}[Divisibility Criterion]
\label{thm:div}
Let $g = a+bi$ with $\gcd(a,b)=1$ and $a^2+b^2>0$. Then $g \mid n$ in $\mathbb{Z}[i]$ if and only if $(a^2+b^2) \mid n$.
\end{theorem}

\begin{proof}
$n/g = n(a-bi)/(a^2+b^2) \in \mathbb{Z}[i]$ iff $(a^2+b^2) \mid na$ and $(a^2+b^2) \mid nb$. Since $\gcd(a,b)=1$, we have $\gcd(a, a^2+b^2) = \gcd(a, b^2) = 1$ and similarly $\gcd(b, a^2+b^2) = 1$, so both conditions reduce to $(a^2+b^2) \mid n$.
\end{proof}

\begin{theorem}[General Divisibility]
For $g = da'+db'i$ with $d = \gcd(a,b)$, $\gcd(a',b')=1$: $g \mid n$ iff $d(a'^2+b'^2) \mid n$.
\end{theorem}

\begin{proof}
$(a'+b'i) \mid (n/d)$ requires $d \mid n$ and $(a'^2+b'^2) \mid (n/d)$, i.e., $d(a'^2+b'^2) \mid n$.
\end{proof}

\begin{theorem}[Decomposition of $R(n)$]
\[
R(n) = \sigma(n) + 2 \sum_{\substack{a' \ge 1,\, b' \ge 1 \\ \gcd(a',b')=1}} \;\sum_{\substack{d \ge 1 \\ d(a'^2+b'^2) \mid n}} d \cdot a'
\]
\end{theorem}

\begin{proof}
Real divisors $d > 0$ contribute $\sigma(n)$. Purely imaginary divisors contribute~$0$. Complex divisors $a+bi$ with $a>0$ come in conjugate pairs $(a+bi, a-bi)$, each contributing real part~$a$. Parametrizing $a = da'$, $|b| = db'$ with $\gcd(a',b')=1$, each pair contributes $2da'$.
\end{proof}

\begin{theorem}[Summation Formula]
\label{thm:sum}
\[
\sum_{n=1}^{N} R(n) = \underbrace{\sum_{d=1}^{N} d\!\left\lfloor \frac{N}{d} \right\rfloor}_{S_1} + 2\!\!\underbrace{\sum_{\substack{a',b' \ge 1,\;\gcd(a',b')=1 \\ a'^2+b'^2 \le N}} \sum_{k=1}^{\lfloor N/(a'^2+b'^2) \rfloor} ka' \!\left\lfloor \frac{N}{k(a'^2+b'^2)} \right\rfloor}_{S_2}
\]
\end{theorem}

\begin{proof}
Exchange the order of summation: for $S_1$, $\sum_n \sigma(n) = \sum_d d\lfloor N/d\rfloor$. For $S_2$, the Gaussian integer $ka'+kb'i$ divides all multiples of $k(a'^2+b'^2)$ up to $N$, of which there are $\lfloor N/(k(a'^2+b'^2))\rfloor$.
\end{proof}

\begin{lemma}[Hyperbola Method]
$S_1$ is computable in $O(\sqrt{N})$ by grouping terms with equal $\lfloor N/d \rfloor$.
\end{lemma}

\begin{proof}
At most $2\sqrt{N}$ distinct quotients exist. Each block of consecutive $d$-values sharing quotient $q$ contributes $q \cdot \sum_{d=d_{\min}}^{d_{\max}} d$, an arithmetic series evaluated in $O(1)$.
\end{proof}

\section{Algorithm}

\begin{verbatim}
S1 = sum_{d=1}^{N} d*floor(N/d)  [hyperbola method, O(sqrt(N))]
S2 = 0
for each coprime (a',b') with a'^2+b'^2 <= N:
    for multiplier k with k*(a'^2+b'^2) <= N:
        S2 += k*a' * floor(N / (k*(a'^2+b'^2)))
    [inner sum via hyperbola method]
answer = S1 + 2*S2
\end{verbatim}

\section{Complexity Analysis}

\begin{itemize}
    \item \textbf{Time:} $O(\sqrt{N})$ for $S_1$; $O(N)$ for $S_2$ (iterating $\sim\!\frac{6}{\pi^2}\cdot\frac{\pi N}{4}$ coprime lattice points). Total: $O(N)$.
    \item \textbf{Space:} $O(1)$.
\end{itemize}

\section{Answer}

\[
\boxed{17971254122360635}
\]

\end{document}
